%%%%%%%%%%%%%%%%%%%%%%%%%%%%%%%%%%%%%%%%%%%%%%%%%%%%%%%%%%%%%%%%%%%%%%%%%%%%%%%%
%% 03_Arquitetura_Sistema_NextGIS.tex
%% Arquitetura de Software e Dados da Plataforma NextGIS Web e PostGIS
%%%%%%%%%%%%%%%%%%%%%%%%%%%%%%%%%%%%%%%%%%%%%%%%%%%%%%%%%%%%%%%%%%%%%%%%%%%%%%%%

\chapter{Arquitetura do Sistema e Plataforma Geoespacial}
\label{cap:arquitetura_sistema}

\section{Justificativa da Plataforma NextGIS Web}

A plataforma NextGIS Web foi selecionada para o grupo de pesquisa em Modelagem Hidrológica por disponibilizar um servidor cartográfico integrado de código aberto, com suporte nativo aos padrões do \textit{Open Geospatial Consortium} (OGC)\footnote{NextGIS Inc. NextGIS Web Documentation and Deployment Guide, 2026.}.

A ferramenta gera serviços web de mapas dinâmicos (\textit{Web Map Service} - WMS), serviços de feições vetoriais (\textit{Web Feature Service} - WFS) e mosaicos de imagens em camadas (\textit{Tile Map Service} - TMS). Os pesquisadores utilizam essas camadas para sobrepor dados de vazão de bacias hidrográficas, redes de microdrenagem pluvial e pontos de descarte de resíduos sobre bases cartográficas de referência.

\section{Decisões de Arquitetura de Software e Dados}

A concepção da solução fundamenta-se em três decisões arquiteturais estruturantes:

\begin{enumerate}
    \item \textbf{Desacoplamento em Duas Camadas de Servidores}: A arquitetura física e lógica separa o processamento em dois servidores independentes. O Servidor de Aplicação Web hospeda o núcleo do NextGIS Web, o proxy reverso Nginx e os módulos de cache. O Servidor de Banco de Dados hospeda o PostgreSQL com PostGIS\footnote{PostGIS Project Steering Committee. PostGIS 3.4 Spatial Database Manual, 2026.} e o armazenamento persistente. A separação garante que consultas espaciais pesadas no banco de dados não afetem a estabilidade e o tempo de resposta da interface web.
    \item \textbf{Isolamento de Dados Relacionais e Matriciais}: Os dados geométricos vetoriais (pontos, linhas e polígonos) são armazenados em tabelas espaciais no PostGIS indexadas por árvores R (GIST). Os arquivos matriciais volumosos, como Modelos Digitais de Elevação (MDE), são mantidos diretamente no sistema de arquivos do servidor com acesso referenciado via catálogo espacial.
    \item \textbf{Ambiente de Treinamento e Homologação em Nuvem}: Enquanto o projeto físico prevê a montagem definitiva da rede no campus da UFS, uma instância virtualizada na nuvem AWS EC2 é empregada para validar a execução dos contêineres e permitir o treinamento imediato dos pesquisadores.
\end{enumerate}

A Figura~\ref{fig:arquitetura_conteineres} apresenta o diagrama da solução NextGIS Web e o fluxo de dados entre os componentes.

\incluirdiagrama{Imagens/arquitetura_conteineres.png}{Arquitetura da solução NextGIS Web, componentes e fluxo de dados geoespaciais.}{fig:arquitetura_conteineres}

\section{Matriz de Aplicações e Dependências}

A Tabela~\ref{tab:aplicacoes_dependencias} detalha os serviços em execução, suas portas de comunicação, dependências diretas e os acordos de nível de serviço (SLA) projetados.

\begin{table}[htbp]
\centering
\small
\caption{Matriz de aplicações, dependências operacionais e acordos de nível de serviço.}
\label{tab:aplicacoes_dependencias}
\begin{tabularx}{\textwidth}{|l|Y|l|l|c|}
\hline
\rowcolor{tableheaderbeige}
\textbf{Serviço} & \textbf{Função no Sistema} & \textbf{Dependências} & \textbf{Portas} & \textbf{SLA} \\ \hline
Nginx Proxy & Terminação HTTP/HTTPS e cache & NextGIS Web & 80, 443 TCP & 99,5\% \\ \hline
NextGIS Core & Servidor de mapas e API OGC & PostGIS, Redis & 8080 TCP (int) & 99,0\% \\ \hline
PostgreSQL/PostGIS & Banco de dados geográfico & Volume persistente & 5432 TCP (int) & 99,9\% \\ \hline
Redis Cache & Fila de tarefas assíncronas & Memória RAM & 6379 TCP (int) & 99,0\% \\ \hline
\end{tabularx}
\end{table}

\section{Implantação da Solução NextGIS Web com Docker Compose}

A padronização dos ambientes de teste e de produção é garantida através de contêineres Docker gerenciados por arquivo \texttt{docker-compose.yml}.

A Listagem~\ref{lst:docker_compose} apresenta a especificação técnica dos serviços. O código foi configurado com redes pontes isoladas e volumes persistentes, evitando sobreposição de portas públicas e quebras de paginação.

\begin{lstlisting}[language=bash,caption={Arquivo de configuração docker-compose.yml para a solução NextGIS Web.},label={lst:docker_compose},basicstyle=\ttfamily\scriptsize]
version: '3.8'

services:
  postgis_db:
    image: postgis/postgis:16-3.4
    container_name: ngw_postgis
    restart: always
    environment:
      POSTGRES_DB: nextgisweb_db
      POSTGRES_USER: nextgis_admin
      POSTGRES_PASSWORD: SecretPostgisPassword2026!
      PGDATA: /var/lib/postgresql/data/pgdata
    volumes:
      - postgis_data:/var/lib/postgresql/data
      - ./backups:/backups
    networks:
      - internal_gis_net

  redis_cache:
    image: redis:7.2-alpine
    container_name: ngw_redis
    restart: always
    networks:
      - internal_gis_net

  nextgisweb:
    image: nextgis/nextgisweb:latest
    container_name: ngw_core
    restart: always
    depends_on:
      - postgis_db
      - redis_cache
    environment:
      NGW_DB_HOST: postgis_db
      NGW_DB_PORT: 5432
      NGW_DB_NAME: nextgisweb_db
      NGW_DB_USER: nextgis_admin
      NGW_DB_PASSWORD: SecretPostgisPassword2026!
      NGW_REDIS_URL: redis://redis_cache:6379/0
    volumes:
      - nextgis_filestore:/opt/nextgisweb/filestore
    networks:
      - internal_gis_net
      - dmz_net

  nginx_proxy:
    image: nginx:alpine
    container_name: ngw_proxy
    restart: always
    ports:
      - "80:80"
      - "443:443"
    depends_on:
      - nextgisweb
    volumes:
      - ./nginx.conf:/etc/nginx/nginx.conf:ro
    networks:
      - dmz_net

volumes:
  postgis_data:
    driver: local
  nextgis_filestore:
    driver: local

networks:
  internal_gis_net:
    driver: bridge
  dmz_net:
    driver: bridge
\end{lstlisting}

\section{Estratégia de Treinamento de Usuários e Homologação em Nuvem}

Para viabilizar o treinamento operacional dos pesquisadores antes da entrega das obras físicas, a aplicação é executada em uma instância virtualizada Amazon EC2 com o sistema Ubuntu Server 24.04 LTS.

A configuração de rede da instância utiliza uma Virtual Private Cloud (VPC) com regras restritivas em seus Grupos de Segurança (\textit{Security Groups}):
\begin{itemize}
    \item Portas TCP 80 e 443 liberadas para acesso público, permitindo que pesquisadores e alunos acessem o portal web e consumam camadas cartográficas;
    \item Porta TCP 22 restrita aos endereços IP de gerência da equipe técnica para manutenção e suporte;
    \item A porta TCP 5432 do PostgreSQL não possui liberação perimetral, permanecendo confinada à comunicação interna entre contêineres na rede Docker.
\end{itemize}
